\documentclass[11pt]{article}

\usepackage[T1]{fontenc}
\usepackage[margin=1in]{geometry}
\usepackage{amsmath, amssymb, amsthm}
\usepackage{enumitem}
\usepackage{hyperref}
\usepackage{booktabs}
\usepackage{listings}
\usepackage{xcolor}

\lstset{
  basicstyle=\ttfamily\small,
  frame=single,
  breaklines=true,
  columns=fullflexible,
  mathescape=true
}

\theoremstyle{definition}
\newtheorem{theorem}{Theorem}[section]
\newtheorem{definition}[theorem]{Definition}
\newtheorem{example}[theorem]{Example}

\title{CS 250 --- Intermezzo: Inference Rules as Trees}
\author{}
\date{}

\begin{document}
\maketitle

\section{A better notation}

In Module 3 we wrote our inference rules in a semi-formal style that looked like little programs---premises on numbered lines, a conclusion at the bottom. That format is fine for getting started, and it has the nice property that it looks like code. But there's a standard notation used in logic textbooks, type theory, and programming language theory that's compact, composable, and worth learning: \emph{Gentzen-style}\index{Gentzen style} inference rules, sometimes called ``natural deduction''\index{natural deduction} style after the proof system Gerhard Gentzen introduced in the 1930s.\footnote{Gentzen was a remarkable logician who also invented the sequent calculus, proved the consistency of Peano arithmetic, and basically invented modern proof theory. All before the age of 30.}

The idea is simple. You draw a horizontal line. Above the line go the \emph{premises}---the things you need to know. Below the line goes the \emph{conclusion}---the thing you get to conclude. Off to the right of the line you write the name of the rule. That's it.

\section{The rules, redux}

Let's revisit the rules from Module 3, but now in this new notation. We'll use \LaTeX's \verb|\dfrac| command, so these are literally just fractions where the numerator is what you know and the denominator is what you conclude.

\subsection{Conjunction}

\noindent\textbf{Introduction:}
If you know $p$ and you know $q$, you can conclude $p \wedge q$:
\[
\dfrac{p \quad q}{p \wedge q}\;\wedge\text{I}
\]
This is the same rule as the \texttt{lstlisting} version from Module 3---premises (a) $p$ and (b) $q$, conclusion $p \wedge q$---just written more compactly.

\noindent\textbf{Elimination:}
If you know $p \wedge q$, you can pull out either piece:
\[
\dfrac{p \wedge q}{p}\;\wedge\text{E}_1
\qquad\qquad
\dfrac{p \wedge q}{q}\;\wedge\text{E}_2
\]

\subsection{Disjunction}

\noindent\textbf{Introduction:}
If you know one disjunct, you can conclude the disjunction:
\[
\dfrac{p}{p \vee q}\;\vee\text{I}_1
\qquad\qquad
\dfrac{q}{p \vee q}\;\vee\text{I}_2
\]
Notice that you get to ``make up'' the other disjunct---if you know $p$, you can conclude $p \vee q$ for \emph{any} $q$. Remember from Module 3 that $p \to (p \vee q)$ really is a tautology.

\noindent\textbf{Elimination (proof by cases):}
\[
\dfrac{p \vee q \quad p \to r \quad q \to r}{r}\;\vee\text{E}
\]
This is the rule we worked through with the eight-row truth table. If you know $p \vee q$ and you can get to $r$ from either side, you've got $r$.

\subsection{Implication}

\noindent\textbf{Elimination (modus ponens):}
\[
\dfrac{p \quad p \to q}{q}\;\to\text{E}
\]
Our old friend. If you know $p$ and you know $p \to q$, you get $q$.

\noindent\textbf{Modus tollens:}
\[
\dfrac{p \to q \quad \lnot q}{\lnot p}\;\text{MT}
\]
If you know $p \to q$ and you know $\lnot q$, you can conclude $\lnot p$. This is the ``backdoor'' version of modus ponens: instead of affirming the antecedent, you deny the consequent. Unlike the contrapositive \emph{technique}, modus tollens as a single inference step is constructively valid.

\noindent\textbf{Introduction:}
This is the rule that was ``a little weird'' in Module 3, and the tree notation actually makes it \emph{clearer}. The idea is: if you can derive $q$ while \emph{temporarily assuming} $p$, then you can discharge that assumption and conclude $p \to q$. We write the discharged assumption in square brackets:
\[
\dfrac{\begin{matrix}[p] \\ \vdots \\ q\end{matrix}}{p \to q}\;\to\text{I}
\]
The $[p]$ means ``$p$ was assumed for the sake of this derivation, and now we're done with it.''\index{discharged assumption} The dots indicate that there's a subderivation---possibly several steps---that gets from $p$ to $q$. Once you apply $\to$I, the assumption $p$ is \emph{discharged}: it's no longer something you're assuming, because you've packaged it up inside the implication.

This is exactly the same as the two-layer structure from Module 3 where we had ``has a proof that assuming $p$, shows $q$''---the bracket notation is just the standard way of writing that in tree form.

\subsection{Negation}

\noindent\textbf{Introduction:}
If assuming $p$ lets you derive a contradiction ($\bot$), then you can discharge that assumption and conclude $\lnot p$:
\[
\dfrac{\begin{matrix}[p] \\ \vdots \\ \bot\end{matrix}}{\lnot p}\;\lnot\text{I}
\]
This works exactly like $\to$I: the brackets around $p$ mean that $p$ was a temporary assumption, and it gets discharged when you apply the rule. In fact, in many formulations of logic $\lnot p$ is simply \emph{defined} as $p \to \bot$, so $\lnot$I is literally a special case of $\to$I.

\noindent\textbf{Elimination (double negation):}
\[
\dfrac{\lnot\lnot p}{p}\;\lnot\text{E}
\]
If you know $\lnot\lnot p$, you can conclude $p$. Note that this rule requires the law of excluded middle---it is not valid in constructive logic. In a proof assistant like Agda or Coq, you would need to explicitly invoke a classical axiom to use it. The direction $p \to \lnot\lnot p$ is constructively valid, but eliminating the double negation is the part that requires classical reasoning (see the discussion in Module 3).

\subsection{Constants}

For completeness:
\[
\dfrac{}{T}\;T\text{I}
\qquad\qquad
\dfrac{\bot}{p}\;\bot\text{E}
\]
$T$ can always be concluded (no premises needed---note the empty space above the line). We write $\bot$ (pronounced ``bottom'' or ``absurdity'') for the constant false---this is the same as $F$ from Module 2, just in standard natural deduction notation. From $\bot$ you can conclude literally anything, which is the ``ex falso'' principle.

\section{Composing rules into proof trees}

Here's where this notation really shines. Because the conclusion of one rule is just a formula, you can plug it in as the premise of another rule. The result is a \emph{tree}\index{proof tree} that grows upward from the final conclusion at the root.

Let's redo the transitivity-of-implication proof from Module 3. We want to show:
\[
\text{Given } p \to q \text{ and } q \to r, \text{ prove } p \to r.
\]
Here's the tree:
\[
\dfrac{
  \dfrac{
    \dfrac{[p] \quad p \to q}{q}\;\to\text{E}
    \quad q \to r
  }{r}\;\to\text{E}
}{p \to r}\;\to\text{I}
\]
Read it from the bottom up: we want $p \to r$, so we use $\to$I, which means we need to derive $r$ from the assumption $[p]$. To get $r$ we use modus ponens with $q \to r$, but we need $q$. To get $q$ we use modus ponens with $[p]$ and $p \to q$. Done.

The assumptions $p \to q$ and $q \to r$ are \emph{undischarged}---they sit at the leaves of the tree as open assumptions. The assumption $[p]$ is discharged by the $\to$I at the root.

Compare this to the \texttt{lstlisting} version:
\begin{lstlisting}
(a) p -> q, by assumption
(b) q -> r, by assumption
(c) proof that assuming p, shows r
          subproof: (i) p, by assumption
                    (j) q, by modus ponens with (i) and (a)
                    show r, by modus ponens with (j) and (b)
shows p -> r by introduction of implication
\end{lstlisting}
Same proof, same logic --- just a different way to lay it out on the page. The tree version makes the structure---which things depend on which---immediately visible.

\section{Reading proof trees}

You can read these trees in two directions, and both are useful.

\textbf{Bottom-up} (goal-directed): start at the conclusion and ask ``what rule could produce this?'' That tells you what you need to prove next, and you work your way up to the leaves. This is the way you typically \emph{construct} a proof---you start with what you want and work backwards until everything is justified.

\textbf{Top-down} (forward reasoning): start at the leaves (your assumptions) and follow the rules downward, seeing what you can derive. This is the way you typically \emph{verify} a proof---you check that each step follows from the ones above it.

This dual reading mirrors a distinction you'll keep encountering: the difference between ``working backwards from what you want to prove'' and ``working forwards from what you know.'' If you've ever debugged a program by starting from the error message and tracing back through the code, you were doing goal-directed reasoning. If you've ever traced through a program with a debugger, stepping forward line by line, you were doing forward reasoning. Both are useful; neither is always better.

\section{Why this matters}

If this is your only logic or discrete math course, you can get by just fine with the \texttt{lstlisting} format from Module 3. But if you go on to take a course in compilers, programming language theory, formal verification, or type theory---and as CS majors, there's a decent chance you will---you're going to see Gentzen-style rules \emph{everywhere}. The typing rules for a programming language are written this way. The operational semantics of a language are written this way. The rules of a type checker, a proof assistant, a static analyzer---all trees, all the time. It's also the notation used in most logic textbooks beyond the intro level, so if you ever crack open a graduate-level logic text, you'll be glad you've seen it before.

Think of this intermezzo as a Rosetta Stone: you now know two ways to write the same rules and can translate between them.

\end{document}
